\section{Der Cauchysche Integralsatz}

\subsection*{\S 1. Stammfunktionen}

In Kap. 4 haben wir gesehen: Ist $f(z) = (z-z_0)^n$ mit $n \neq -1$, so verschwindet das Integral längs eines geschlossenen Weges. Für $n=-1$ ist es $2\pi i$. Das erinnert an die reelle Analysis: $x^n$ hat für $n \neq -1$ eine Stammfunktion $\frac{1}{n+1}x^{n+1}$. Für $n=-1$ ist es $\ln|x|$.

\vspace{1em}
\noindent\textbf{Def: 5.1 (Stammfunktion)} \\
Sei $G \subset \mathbb{C}$ offen. Eine Funktion $F: G \rightarrow \mathbb{C}$ heißt \textit{Stammfunktion} einer Funktion $f: G \rightarrow \mathbb{C}$, falls $F$ auf $G$ holomorph ist mit $F'(z) = f(z)$ für alle $z \in G$.

\vspace{1em}
\noindent\textbf{Satz: 5.2 (Hauptsatz für Wegintegrale)} \\
Sei $G \subset \mathbb{C}$ offen, $f: G \rightarrow \mathbb{C}$ stetig und $F: G \rightarrow \mathbb{C}$ eine Stammfunktion von $f$. Ist $\gamma: [a,b] \rightarrow G$ ein Integrationsweg (stückw. $C^1$-Weg), so gilt:
\[
\int_{\gamma} f(z) \, dz = F(\gamma(b)) - F(\gamma(a))
\]
Insbesondere gilt: Ist $\gamma$ geschlossen (d.h. $\gamma(a) = \gamma(b)$), so ist $\int_{\gamma} f(z) \, dz = 0$.

\vspace{1em}
\noindent\textbf{Beweis:} \\
Sei $\gamma$ zunächst ein $C^1$-Weg. Nach Def. des Wegintegrals und der Kettenregel gilt:
\[
\int_{\gamma} f(z) \, dz = \int_{a}^{b} f(\gamma(t)) \cdot \gamma'(t) \, dt = \int_{a}^{b} F'(\gamma(t)) \cdot \gamma'(t) \, dt = \int_{a}^{b} \frac{d}{dt} \big(F(\gamma(t))\big) \, dt
\]
Nach dem HDI für reelle/komplexe Funktionen über reelle Intervalle (Satz 4.8) gilt:
\[
= \big[F(\gamma(t))\big]_{a}^{b} = F(\gamma(b)) - F(\gamma(a))
\]
Ist $\gamma$ nur stückweise $C^1$, so zerlegt man das Intervall $[a,b]$ in Teilintervalle, auf denen $\gamma$ jeweils $C^1$ ist, und erhält per Teleskopsumme das gleiche Ergebnis. \hfill $\square$

\vspace{1em}
\noindent\textbf{Folgerung: 5.3} \\
Die Funktion $f(z) = \frac{1}{z}$ besitzt in $G = \mathbb{C} \setminus \{0\}$ keine Stammfunktion.

\vspace{1em}
\noindent\textbf{Beweis:} \\
Wäre $F$ eine Stammfunktion von $f(z) = \frac{1}{z}$ auf $\mathbb{C} \setminus \{0\}$, so müsste für jeden geschlossenen Weg $\gamma$ in $G$ gelten: $\int_{\gamma} \frac{1}{z} \, dz = 0$. \\
Für den Einheitskreis $\gamma(t) = e^{it}$ mit $t \in [0, 2\pi]$ gilt aber nach Satz 4.11:
\[
\oint_{\gamma} \frac{1}{z} \, dz = 2\pi i \neq 0
\]
Widerspruch! \hfill $\square$

\vspace{1em}
\noindent\textbf{Bemerkung:} \\
Auf dem einfach zusammenhängenden Gebiet $G' = \mathbb{C} \setminus (-\infty, 0]$ (geschlitzte Ebene) besitzt $f(z) = \frac{1}{z}$ jedoch eine Stammfunktion, nämlich den Hauptzweig des komplexen Logarithmus $\log(z)$.

\vspace{1em}
\noindent\textbf{Satz: 5.4 (Existenz von Stammfunktionen)} \\
Für eine stetige Funktion $f: G \rightarrow \mathbb{C}$ auf einem Gebiet $G \subset \mathbb{C}$ sind folgende Aussagen äquivalent:
\begin{enumerate}
    \item[i)] $f$ besitzt eine Stammfunktion $F$ auf $G$.
    \item[ii)] Für je zwei Integrationswege $\gamma_1, \gamma_2$ in $G$ mit gleichem Anfangs- und Endpunkt gilt: $\int_{\gamma_1} f(z) \, dz = \int_{\gamma_2} f(z) \, dz$ (Wegunabhängigkeit).
    \item[iii)] Für jeden geschlossenen Integrationsweg $\gamma$ in $G$ gilt: $\oint_{\gamma} f(z) \, dz = 0$.
\end{enumerate}

\vspace{1em}
\noindent\textbf{Beweis von Satz 5.4:}
\begin{itemize}
    \item[\textbf{i) $\implies$ iii):}] Folgt direkt aus Satz 5.2 (Hauptsatz für Wegintegrale), da $\gamma(a) = \gamma(b)$.
    
    \item[\textbf{iii) $\implies$ ii):}] Seien $\gamma_1, \gamma_2$ zwei Wege von $z_1$ nach $z_2$. Dann ist der zusammengesetzte Weg $\gamma := \gamma_1 \oplus (-\gamma_2)$ ein geschlossener Weg. Nach Voraussetzung gilt:
    \[
    0 = \oint_{\gamma} f(z) \, dz = \int_{\gamma_1} f(z) \, dz + \int_{-\gamma_2} f(z) \, dz = \int_{\gamma_1} f(z) \, dz - \int_{\gamma_2} f(z) \, dz
    \]
    Daraus folgt $\int_{\gamma_1} f(z) \, dz = \int_{\gamma_2} f(z) \, dz$.

    \item[\textbf{ii) $\implies$ i):}] Wähle einen festen Punkt $z_0 \in G$. Da $G$ ein Gebiet (offen und wegzusammenhängend) ist, gibt es zu jedem $z \in G$ einen Weg von $z_0$ nach $z$. Wir definieren:
    \[
    F(z) := \int_{\gamma_z} f(\zeta) \, d\zeta
    \]
    wobei $\gamma_z$ ein beliebiger Weg von $z_0$ nach $z$ in $G$ ist. Aufgrund von ii) ist $F(z)$ wohldefiniert (unabhängig von der Wahl von $\gamma_z$).
    
    Zeige nun, dass $F'(z) = f(z)$: Sei $z \in G$ und $h \in \mathbb{C}$ klein genug, sodass $z+h \in G$ (da $G$ offen). Da die Strecke $[z, z+h]$ ganz in $G$ liegt, können wir den Weg nach $z+h$ wählen als $\gamma_z \oplus [z, z+h]$.
    \[
    F(z+h) - F(z) = \int_{\gamma_z \oplus [z, z+h]} f(\zeta) \, d\zeta - \int_{\gamma_z} f(\zeta) \, d\zeta = \int_{[z, z+h]} f(\zeta) \, d\zeta
    \]
    Es gilt:
    \[
    \frac{F(z+h) - F(z)}{h} - f(z) = \frac{1}{h} \int_{[z, z+h]} \big(f(\zeta) - f(z)\big) \, d\zeta
    \]
    Da $f$ stetig ist, gibt es zu $\epsilon > 0$ ein $\delta > 0$, sodass $|f(\zeta) - f(z)| < \epsilon$ für $|\zeta - z| < \delta$. Für $|h| < \delta$ gilt also mit der Standard-Abschätzung (Satz 4.7):
    \[
    \left|\frac{F(z+h) - F(z)}{h} - f(z)\right| \le \frac{1}{|h|} \cdot \max_{\zeta \in [z, z+h]} |f(\zeta) - f(z)| \cdot \text{Länge}\big([z, z+h]\big) < \frac{1}{|h|} \cdot \epsilon \cdot |h| = \epsilon
    \]
    Für $h \rightarrow 0$ folgt $F'(z) = f(z)$. Also ist $F$ eine Stammfunktion von $f$. \hfill $\square$
\end{itemize}
